\section{Introduction}
\label{sec:intro}

A robot that follows instructions such as ``pick up the pear by its stem'' must find the named object and predict a grasp on it.
Grasp-Anything++~\cite{vuong2024language} frames this as language-driven grasp detection: given an RGB image and a prompt, predict a 2D grasp rectangle $g = (x, y, w, h, \theta)$ for a two-finger gripper.
Recent methods improve the success rate with diffusion models~\cite{vuong2024language,nguyen2024lightweight}, mask-guided attention~\cite{vo2024maskgrasp} or state-space models~\cite{nguyen2024graspmamba}.

Before proposing a method, we ask where a simple baseline fails.
We reproduce GR-ConvNet + CLIP~\cite{kumra2020antipodal,vuong2024language} and classify 1,000 of its failures with two vision-language models.
About 56\% of the failures grasp another object than the prompted one (\cref{tab:errors}); a wrong grasp shape on the right object accounts for only about 29\%.
We also find that the part-level labels of the dataset mostly repeat the object-level ones, so the prompt mainly selects the object.
The baseline's loss only says where to grasp: nothing tells its encoder, trained from scratch, which object a prompt names.

Our contributions are:
(1) a failure analysis of a reproduced baseline on Grasp-Anything++ that identifies grounding as the main error;
(2) a region-text contrastive grounding loss with hard negatives taken from the other objects of the same scene, used in training only;
(3) a gain of 7.5 points on unseen classes over the same baseline trained under identical conditions.
